\section{The scalar error at the trial potential}\label{sec:scalar}

The first error in the capacity comparison is the pairing of the
propagated source with the trial potential itself.  We estimate it
using Lemma~\ref{lem:duhamel}.  Its first term benefits from signed
cancellation on isolated strands.  Its second term contains two
sources, and heat-kernel decay and spatial localization make their
joint contribution smaller.  All expectations below use the actual
pin law $\PP_P$.

\begin{proposition}[Scalar estimate]\label{prop:scalar}
With the scales of \eqref{eq:scales}, under the interiority and pin
assumptions of Section~\ref{sec:profiles},
\begin{equation}\label{eq:scalarclaim}
 \EE_P\bigl[\ind_{\mathcal G}|\ip{\mu_+}{h}|\bigr]
                         =o(\beta^{-5/2}).
\end{equation}
The error is uniform in the finite cube, the deterministic pin set,
and the injection slot.
\end{proposition}

\subsection{A single source paired with the free profile}

For forward time $s$ put
\[
 d=\sqrt{1+s},\qquad d'=\sqrt{1+\beta-s}.
\]
The reverse profile $\bar b_s$ has value bound $C(G+P_{d'})$.
The basic weight integral is
\begin{equation}\label{eq:scalarweightintegral}
 \int_0^\beta\sum_{r_z\le4R}
       (W+W_d)(z)(G+P_{d'})(z)\,ds
                  \le C\beta^{-1}(1+\log\beta).
\end{equation}
To verify it, first sum the smooth product by radial shells:
\[
 \sum_{r_z\le4R}W(z)G(z)
 \le C\beta^{-2}
       \left(1+\int_0^{4R}\frac{r^2\,dr}{(D_*+r)^3}\right)
 \le C\beta^{-2}(1+\log\beta).
\]
For $s\le\beta/2$, Lemma~\ref{lem:freebounds} gives $P_{d'}\le CG$,
and
\[
 \sum_z W_d(z)G(z)
 \le \sup G\sum_z W_d(z)\le C\beta^{-3/2}d^{-1}.
\]
Integration over this half-period costs at most $C\beta^{-1}$.
For $s\ge\beta/2$, use $W_d\le CW$ and
$\sum_zP_{d'}(z)\le C$, so the remaining product costs at most
$C\sup W=C\beta^{-2}$ per unit time.  This proves
\eqref{eq:scalarweightintegral}.

Consider first the blocks contained in
$[\beta^{1/4},\beta-\beta^{1/4}]$.  On each such block
$d,d'\ge\beta^{1/8}$ and their weights are comparable over the block
and over distances $CL$.  At a deterministic initial slice,
Proposition~\ref{prop:snapshot} bounds the probability of an obstacle
at a prescribed slot by $C\lambda$.  If it is not isolated, a second
distinct obstacle lies within $c_2L$, whence the same proposition and
a union bound give
\begin{equation}\label{eq:nonisolatedprobability}
 \PP_P(\text{the prescribed slot is an obstacle and is not isolated})
                       \le CL^3\lambda^2.
\end{equation}
Here and below fixed slot multiplicities are included in constants.

For a nonisolated strand starting over $z$, the total variation
bound in Lemma~\ref{lem:cancellation} bounds its scalar contribution
by
\[
 CL(W+W_d)(z)(G+P_{d'})(z).
\]
For an isolated strand subtract from the testing profile $\bar b$
its value at the strand's initial site and time, at any one slot.
The constant value pairs with the signed source, whose bound is
$CL^{3/4}(W+W_d)(z)$.  The remaining variation of the testing profile
is at most
\begin{equation}\label{eq:profilevariation}
 CL(W+W_{d'})(z)
 \le C L\beta^{-1/8}(G+P_{d'})(z).
\end{equation}
Spatial variation follows from the difference bounds over distance
$CL$; time variation over the block follows from the time derivative
bounds.  The last inequality uses $W\le G/D_*$ and
$W_{d'}=P_{d'}/d'$.

Take expectations, use \eqref{eq:nonisolatedprobability}, and sum the
blocks.  Their lengths are comparable to $L$, so their weight sums
are bounded by $L^{-1}$ times
\eqref{eq:scalarweightintegral}.  The contribution of all these
blocks is at most
\begin{equation}\label{eq:centralfirstterm}
 C\beta^{-1}(1+\log\beta)
   \left[L^3\lambda^2
          +\lambda\bigl(L^{-1/4}+L\beta^{-1/8}\bigr)\right]
                         =o(\beta^{-5/2}).
\end{equation}
The three powers of $\beta$ are respectively
$-4+3\delta$, $-5/2-\delta/4$, and $-21/8+\delta$.
The middle term is the one for which signed cancellation is needed.

The remaining blocks lie within $2\beta^{1/4}$ of an endpoint, for
large $\beta$, because $L\ll\beta^{1/4}$.  Use the bin total
variation bound from Lemma~\ref{lem:sourcebins}, without cancellation.
Weights at times and sites within a bin change by at most $H_\beta$,
so the expected contribution is bounded by $H_\beta\lambda$ times
the integral in \eqref{eq:scalarweightintegral} restricted to these
endpoint intervals.  On the first interval the preceding estimates
give
\[
 C\left[\beta^{1/4}\beta^{-2}(1+\log\beta)
       +\beta^{-3/2}\int_0^{2\beta^{1/4}}
                              (1+s)^{-1/2}\,ds\right]
 \le C\left[\beta^{-7/4}(1+\log\beta)+\beta^{-11/8}\right].
\]
The last interval costs at most
$C\beta^{-7/4}(1+\log\beta)$.  Multiplication by
$H_\beta\lambda$ gives powers at most $-23/8$ and $-13/4$,
both strictly below $-5/2$.

The cutoff background is estimated separately.  On its annulus the
reverse profile has size at most $C/(\beta R)$, so
\begin{equation}\label{eq:scalarbackground}
 \left|\int\ip{\gamma_s|_{H_s}}{\bar b_s}\,ds\right|
 \le C\beta R^3 b_*\frac1{\beta R}
 =\frac C{\beta R}=o(\beta^{-5/2}).
\end{equation}
We have therefore bounded the first term of
\eqref{eq:doubledduhamel}.  It remains to control the interaction
between its two source collections.

\subsection{Snapshot moments for spatial boxes}

The source locations at two different times are correlated.  The
following consequence of the snapshot estimate controls their
pairing without an independence assumption.

\begin{lemma}[Box moment bound]\label{lem:boxmoment}
Partition the sites into cubes $Q$ of side comparable to $F\ge1$.
At any deterministic slice let $\mathcal O$ be its obstacle slots,
and let $w$ be a nonnegative deterministic site weight of finite
support.  With slot multiplicity in the sums, set
$A_Q=\sum_{(z,a)\in\mathcal O,\ z\in Q}w(z)$.  Then
\begin{equation}\label{eq:boxmoment}
 \left(\EE_P\sum_Q A_Q^2\right)^{1/2}
 \le C\left[
  \sqrt\lambda\,\|w\|_2
  +\lambda\left\{\sum_Q
                     \left(\sum_{z\in Q}w(z)\right)^2\right\}^{1/2}
                 \right].
\end{equation}
If two such sums are taken at different deterministic slices, the
expected sum of products over pairs of anchors at distance at most
$CF$ is bounded by a constant times the product of their respective
right sides in \eqref{eq:boxmoment}.
\end{lemma}

\begin{proof}
Expand $\sum_Q A_Q^2$.  Terms using the same slot are bounded by
$C\lambda\|w\|_2^2$ by Proposition~\ref{prop:snapshot}.  Terms using
two distinct slots, including two slots at the same site, are bounded
by $C\lambda^2\sum_Q(\sum_{z\in Q}w(z))^2$.  Taking a square root
proves \eqref{eq:boxmoment}.
At a seam the same bound holds on either side: each $A_Q$ depends
only on the number of obstacle slots at each site, and these numbers
are preserved by the seam.

Distance at most $CF$ restricts the two anchors to cubes whose
indices differ by one of a bounded number of offsets.  For each
offset, Cauchy--Schwarz in the product of the picture probability
space and the cube index set gives
\[
 \EE_P\sum_Q A_Q A'_{Q+v}
 \le\left(\EE_P\sum_Q A_Q^2\right)^{1/2}
      \left(\EE_P\sum_Q (A'_Q)^2\right)^{1/2}.
\]
Sum the bounded number of offsets.  Both factors use their own
snapshot estimate, so no joint law at the two times is required.
\end{proof}

In the applications, the weights are restricted to the deterministic
anchor region $B_{4R}(x)$.  We bound their norms by the corresponding
unrestricted envelopes.  For these weights, denote the resulting bounds by
\begin{equation}\label{eq:JB}
 \begin{aligned}
 J(d)&=C\bigl(\beta^{-3/4}d^{-5/2}
                         +\beta^{-3/2}d^{-1}\bigr),
       &&w=W_d,\\
 B(F)&=C\bigl(\beta^{-2}+\beta^{-5/2}F\bigr),
       &&w=W.
 \end{aligned}
\end{equation}
To check the first line, direct shell sums give
$\|W_d\|_2\le Cd^{-5/2}$ and $\|W_d\|_1\le Cd^{-1}$;
the square root of the squared box sums is no greater than the
total weight.  For the second line,
$\|W\|_2\le C\beta^{-1}D_*^{-1/2}$.
It remains to bound the squared box sums by $C\beta^{-2}F^2$.
Use the larger envelope $C\beta^{-1}(1+r_z)^{-2}$.
The bounded number of cubes within distance $CF$ of $x$ each have
weight sum at most $CF/\beta$.  At distance comparable to $nF$,
each cube has weight at most $C\beta^{-1}F n^{-2}$ and there are
$O(n^2)$ cubes in such a shell.  Summing their squares over $n\ge1$
gives $C\beta^{-2}F^2\sum_{n\ge1}n^{-2}$, as claimed.

The time sum of $J$ over the bins of either traversal satisfies
\begin{equation}\label{eq:Jsum}
 \sum_{\mathrm{bins}}J\bigl(\sqrt{1+s_{\mathrm{bin}}}\bigr)
                           \le C\beta^{-3/4}.
\end{equation}
Indeed the bin lengths are bounded above and below, so
$\sum(1+s_{\mathrm{bin}})^{-5/4}\le C$ and
$\sum(1+s_{\mathrm{bin}})^{-1/2}\le C\sqrt\beta$.
The latter gives $C\beta^{-1}$, which is smaller than the
right side of \eqref{eq:Jsum}.

\subsection{Two sources separated in time}

First remove interactions involving cutoff background.  Counting
mass and supremum contractions of the actual kernels imply that
two background sources cost at most
\[
 C\beta^2b_*^2R^3=CR^{-3}.
\]
For one background source, integrate its time variable and use its
supremum bound.  The resulting cost is at most $C\beta b_*$ times
the total variation of the other, obstacle, source collection.
Lemmas~\ref{lem:sourcebins} and \ref{lem:boxmoment}, or simply the
one-point snapshot bound, give
\begin{equation}\label{eq:expectedsourcevariation}
 \EE_P\left[\ind_{\mathcal G}
                \int\|s^{\pm}_{\mathrm{obs}}\|_1\right]
                  \le H_\beta\lambda(R+\sqrt\beta).
\end{equation}
For the smooth weight, use
$\sum_{r_z\le4R}W(z)\le C\beta^{-1}R$ and $O(\beta)$ bins.
For the point weight use $\sum_zW_d(z)\le C/d$ and sum in time.
Thus all terms with background are $o(\beta^{-5/2})$.

We can now assume both sources are obstacle sources.  Group their
bin pairs by dyadic scales $T\in\{1,2,4,\ldots\}$, up to a constant
times $\beta$, according to $1+t-s\asymp T$.  The scale $T=1$
also includes a bounded number of neighboring and identical bins.
Equivalently one may group by bin indices; the bin length bounds
preserve comparability away from this first scale.  For a fixed bin
there are $O(T)$ partner bins at scale $T$.

First replace each intervening propagator by its version killed on
exiting $B_{R_*}(x)$.  Every source starts in $B_{R+2}(x)$, so
Lemma~\ref{lem:localization} bounds the discarded mass by a rapidly
small quantity, uniformly on $\mathcal G$.  The source supports and
total variations are bounded by fixed powers of $\beta$, by
\eqref{eq:backgroundbound} and Lemma~\ref{lem:sourcebins}; hence the
aggregate loss in the double pairing is still rapidly small.
Lemma~\ref{lem:heatbound} now bounds each entry of the killed
kernel by $CT^{-3/2}$.  The displacement assertion in
Lemma~\ref{lem:localization} further permits us, with another rapidly
small aggregate error, to retain only anchor sites at distance at most
\begin{equation}\label{eq:scalarinteractionradius}
 F=F(T):=C(1+T/L)L^{3/4}.
\end{equation}
The $C\log\beta$ shifts from the bin anchors to insertion sites
are absorbed by $L^{3/4}$.

Apply Lemma~\ref{lem:boxmoment} with boxes of side comparable to
$F(T)$.  The two bin total variations contribute only $H_\beta$
times their anchor weights.  We treat the three possible weight
products in turn.

\paragraph{Two smooth weights.}
There are $O(\beta)$ first bins and $O(T)$ partners per first bin.
Their contribution at scale $T$ is at most
\begin{equation}\label{eq:WWscale}
 H_\beta\beta T^{-1/2}B(F)^2
 \le H_\beta\left[\beta^{-3}T^{-1/2}
                        +\beta^{-4}T^{-1/2}F^2\right].
\end{equation}
Since $F^2\le C(L^{3/2}+T^2L^{-1/2})$, the dyadic sum is at most
\begin{equation}\label{eq:WWsum}
 H_\beta\left[
       \beta^{-3}+\beta^{-4}L^{3/2}
                       +\beta^{-5/2}L^{-1/2}\right]
                         =o(\beta^{-5/2}).
\end{equation}
Here $\sum_TT^{-1/2}\le C$ and
$\sum_{T\le C\beta}T^{3/2}\le C\beta^{3/2}$.
The final term explains why a displacement bound smaller than
$T$ is needed in \eqref{eq:scalarinteractionradius}.

\paragraph{One smooth and one point weight.}
Sum the point weight over its own elapsed-time bins using
\eqref{eq:Jsum}; each such bin has $O(T)$ smooth partners.
The scale-$T$ bound is
\[
 H_\beta T^{-1/2}B(F)\beta^{-3/4}.
\]
Inserting \eqref{eq:JB} and \eqref{eq:scalarinteractionradius},
and summing dyadically, gives at most
\begin{equation}\label{eq:WPsum}
 H_\beta\left[
  \beta^{-11/4}+\beta^{-13/4}L^{3/4}
                         +\beta^{-11/4}L^{-1/4}\right]
                           =o(\beta^{-5/2}).
\end{equation}
This time the growing sum is
$\sum_{T\le C\beta}T^{1/2}\le C\sqrt\beta$.
The calculation applies whichever direction supplies the point
weight, because \eqref{eq:Jsum} holds in either elapsed-time order.

\paragraph{Two point weights.}
When $t-s\le\beta/2$, their elapsed profile times satisfy
\[
 s+(\beta-t)=\beta-(t-s)\ge\beta/2.
\]
At least one is therefore at least $\beta/4$, and its point weight
is bounded by $CW$ by Lemma~\ref{lem:freebounds}.  These pairs are
already covered by \eqref{eq:WPsum}.  The remaining pairs have
separation comparable to $\beta$.  Use the entry bound
$C\beta^{-3/2}$ and the unrestricted two snapshot pairing from
Lemma~\ref{lem:boxmoment}, with a single box containing the source
support.  Summing both time variables by \eqref{eq:Jsum} costs at most
\begin{equation}\label{eq:PPsum}
 H_\beta\beta^{-3/2}
       \left(\sum_{\mathrm{bins}}J(d)\right)^2
                   \le H_\beta\beta^{-3}.
\end{equation}

This finishes the estimate of the second term in
\eqref{eq:doubledduhamel}.  To conclude, the smallest power gains
in the entire scalar calculation are $L^{-1/4}=\beta^{-\delta/4}$
in \eqref{eq:centralfirstterm}, $\beta^{-\delta/4}$ in
\eqref{eq:scalarbackground}, and $L^{-1/2}=\beta^{-\delta/2}$ in
\eqref{eq:WWsum}.  Each dominates every fixed logarithmic power
because $\delta=1/200>0$.  All other displayed terms have larger
power margins.  Taking the expected absolute value in
Lemma~\ref{lem:duhamel} proves Proposition~\ref{prop:scalar}.
